%!TEX program = ksRunCluttex
%! KSRC --lualatex --synctex=1 --max-iterations=2
\documentclass{simplman}

\usepackage{inlinegraphicx}
\usepackage{ninecolors}
\usepackage[shortverb=false,private,backcolor=gray9!30,framecolor=gray5!50]{srcandtext}

\newfontfamily\exemfont{Noto Sans CJK KR Light}[Script=Hangul,Language=Korean,Renderer=OpenType,CompressPunctuations]

\usepackage{microtype}
\usepackage{wasysym}
\usepackage{kswdash}
\renewcommand*{\KSWdashCode}{FF5E}
\usepackage[localenv,parentype=moreparens]{kskorsymb}
\let\ksemph\respectdotemph
\let\ksps\ksparensmall
\usepackage{ksvqt}
\newcommand*\pmime{\<3\textsf{pm IME GUI}>}

\newcommand*\env[1]{\<2\textit{#1}>}
\let\pkg\stp
\let\opt\stopt
\let\file\path
\def\util#1{\texttt{#1}}

\newcommand*\guielem[1]{\colorbox{blue8!20}{\sffamily[#1]}}

\ifxetex
	\def\flat{\symbol{"266D}}
	\def\natural{\symbol{"266E}}
\fi

\ExplSyntaxOn

\NewDocumentEnvironment {widetext} {}
{
	\begin{adjustwidth}{-.5\spinemargin}{0pt}
}
{
	\end{adjustwidth}
}

\NewDocumentCommand \applesymbolfont { m }
{
	\fontspec{Apple~Symbols.ttf}[Scale=#1]
}

\ExplSyntaxOff

\usepackage{multirow}

\maxtocdepth{subsection}

\setlength\cftsectionindent{2.2em}
\setlength\cftsubsectionindent{4.2em}
\counterwithout{figure}{chapter}
\counterwithout{table}{chapter}

\usepackage{printlen}
\usepackage{tabularray}

\def\OK{○}
\def\SEM{△}
\def\NON{×}
\def\textcaret{\char"005E}

\ifxetex
\usepackage{ksruby}
\fi

\begin{document}

\begin{titlingpage}
\calccentering{\unitlength}
\begin{adjustwidth}{\unitlength}{-\unitlength}
\begin{center}
\inlinegraphics[scale=4]{images/pmimegui.png}
\huge \pmime\ 활용 지침서
%\inlinegraphics[scale=2]{mtxvisualhelper.png}
\\
\large Nova de Hi\quad\smaller June--August 2026, ver 0.13.0
\end{center}

\vspace{-\baselineskip}

\aliaspagestyle{chapter}{empty}
\tableofcontents*
\end{adjustwidth}

\end{titlingpage}

\spaceskip=.4em plus .12em minus .05em

\chapter{목적과 용도}

\addhangulfontfeature{Script=Hangul}

\pmime 는 문서를 작성할 때 사용하는 입력 보조 프로그램이다. 
한국어 \LaTeX\ 작업 중이라도 본문 내용에 따라 옛한글, 일본 문자, 일본식 한자, 중국어 간자체 등의 입력이 필요할 때가 있다. 그러나 매번 시스템 입력기를 \ksps{새로 설치하거나} 전환하는 방식은 매끄러운 집필 흐름을 방해한다. 특히 기본 입력기 외에 다수의 추가 입력기가 설치되어 있으면, 이들 사이를 전환하는 것만으로도 사용자의 피로도가 크게 높아진다. 만약 필요한 소량의 텍스트를 간편하게 생성하고, 이를 \ksps{클립보드를 통하여} 작성 중인 문서 소스에 즉시 반영할 수 있다면 편집 효율을 극대화할 수 있다. \pmime 는 바로 이러한 불편을 해소하기 위해 고안되었다.

필자가 최근 작업 중인 문헌은 1920\~1930년대 문화와 관련된 것으로, 이 시기의 텍스트는 이른바 `1933년 이전 표기법'으로 되어 있다. 필자는 이를 `20세기 초반 표기법'이라 부르고 있으나, 기술적으로 더 정확한 표현은 `유니코드 한글 음절 문자 영역을 벗어나는 한글 표기'일 것이다. 이는 15세기 이래의 이른바 `옛한글'과는 성격이 다소 다르므로 구별하여 취급하는 것이 타당하다. 기술적으로 유니코드 자모 한글을 사용하여 입출력한다는 점은 동일하지만 언어는 현대어와 거의 같아서 현대 표기로 전사할 수 있으나 당시 표기가 지금과 다르다는 의미에서 그러하다.

이 표기법의 가장 큰 특징은 된소리를 표기할 때 각자병서(겹자음)를 쓰지 않고 `ㅅ\nobreak 계 합용병서'로 경음을 표기한다는 점이다. 예: 가실ᄯᅢ에는、 ᄭᅩᆺ츨. 이러한 표기들을 입력하는 과정에서 겪은 어려움을 해결하고자, 자모 단위의 한글 입력 방식을 지원하는 \href{https://ctan.org/pkg/pmhanguljamo}{\pkg{pmhanguljamo}} 패키지의 메커니즘을 차용하여 \textsf{pmhangul GUI}라는 애플리케이션을 제작하였다. 

이에 더하여, 한글 입력 방식을 확장함과 동시에 20세기 초반에 사용되던 문장부호 몇 가지를 추가하였다. 또한, 시스템 입력기를 전환하지 않고도 일본 문자를 비롯하여 일본식 및 중국식 한자를 간편하게 생성할 수 있도록 구현한 결과물이 현재의 \pmime 이다.

\pmime 는 처음부터 시스템 입력기를 대체할 목적으로 개발된 것이 아니다. 비록 \LaTeX{} 작업 환경을 최우선으로 염두에 두고 있으나, 필요에 따라 다른 텍스트 에디터나 워드프로세서에서도 앞서 언급한 문자들을 효율적으로 생성하고 입력하는 데 유용한 도움을 줄 수 있을 것이다.

\chapter{설치 및 실행}

설치 및 실행에 필요한 프로그램의 다운로드 링크는 \href{https://nonametex.imweb.me/?bmode=view&idx=171896425}{\<2{\ttfamily nonametex}:171896425>} 게시글에
제시되어 있다.

이 앱을 사용하려면 Noto Sans CJK KR 폰트가 시스템에 설치되어 있어야 한다.

\section{Windows}

인스톨러 프로그램 \file{pmIMEGUI-Installer.exe}를 실행한다. 설치 시에 바탕 화면에
아이콘을 둘 것인지 여부를 선택할 수 있다.
\begin{itemize}\tightlist
\item \file{pm IME GUI.exe}
\item \file{pmimec.exe}
\end{itemize}
두 개의 실행 파일이 설치되며 이 가운데 바탕화면의 아이콘을 클릭하여 실행되는 것은 \file{pm IME GUI.exe}이다. \file{pmimec.exe}는 명령행 유틸리티로서 그 사용법에 대해서는 \pageref{sec:cli}페이지의 \<!\titleref{sec:cli}>장에서 설명한다.

\section{macOS}

인스톨러 \file{pmIMEGUI.dmg} 또는 \file{pmIMEGUI_x64.dmg}를 다운로드하여 마운트하면
app 번들과 \file{json.zip} 두 개의 파일이 있다. \verb|_x64|는 인텔맥용이다.
이 가운데 app 번들을 \file{/Applications/}로 복사하고 \file{json.zip}은 적당한 곳에 복사하여 압축을 해제한다.

\begin{enumerate}[\quad (1)]\oblivoirlist
\item app 번들에 대하여 quarantine을 한번 해제해주어야 한다. 이 프로그램이 서명되어
있지 않기 때문에 그대로 실행하면 아마도 휴지통으로 옮기려고 할 것이다.
\begin{verbatim}
xattr -d com.apple.quarantine /Applications/pm\ IME\ GUI.app
\end{verbatim}
\item \file{json.zip}을 풀어서 나오는 네 개의 \file{*.json}을 다음 위치로 복사한다.
\begin{verbatim}
mkdir -p ~/.config/pmimegui
cp *.json ~/.config/pmimegui/
\end{verbatim}
\end{enumerate}

\section{Python 스크립트 실행}

Linux 등에서는 별도의 실행 프로그램이 제공되지 않으므로 Python 스크립트를 직접
실행한다.

이 프로그램은 Python과 PyQt6로 작성되었다. CLI(명령행 인터페이스)만을 사용하려면
별도의 추가 설치 없이 python을 바로 실행할 수 있으나 GUI 실행을 위해서는
다음 두 가지 패키지가 필요하다.
\begin{enumerate}\tightlist
\item PyQt6
\item pyautogui
\end{enumerate}
Python은 3.12 이후 버전을 요구한다. 만약 가상환경에서 실행해보려 한다면, 
\begin{verbatim}
python -m venv myvenv
source myvenv/bin/activate
\end{verbatim}
가상환경을 활성화한 다음 \util{pip}로 두 개의 패키지를 설치한다.
\begin{verbatim}
pip install PyQt6 pyautogui
\end{verbatim}
그리고 다음과 같이 한다.
\begin{verbatim}
python pmimegui.py
\end{verbatim}
이 스크립트는 인자가 주어지지 않으면 GUI를 띄운다. 반면 하나 이상의 인자가 들어오면
명령행 인터페이스로 동작한다. 명령행 인터페이스를 테스트해보려면 인자를 부여하라.
\begin{verbatim}
python pmimegui.py -h
\end{verbatim}

Linux 운영체제라면 위의 macOS에서 설명한 바와 같이 \file{*.json}을 복사해두는
것이 좋다.
\begin{verbatim}
mkdir -p ~/.config/pmimegui
cp *.json ~/.config/pmimegui/
\end{verbatim}


\chapter{기본 인터페이스}

\begin{widetext}
\raggedleft
\includegraphics[scale=.45,raise=.8cm]{scr01}
\includegraphics[scale=.33]{scr01-1}
\end{widetext}

Windows와 macOS에서 실행한 그림이다. 이후 설명을 위한 스크린샷은 macOS의 것을
보일 테지만 기능은 동일하다.

이 프로그램은 어디까지나 \ksemph{보조} 프로그램이므로 창이 너무 커져서 작업 영역을
가려서는 안된다. 창의 크기는 오른쪽 아래의 모서리를 끌어서 높이를 조절할 수 있고
입력 영역과 출력 영역의 비율도 가운데 스플릿 바를 끌어서 조절할 수 있다.
이전 프로그램에 있던 ``미니모드''는 없어졌지만 기능이 너무 복잡해져서 미니모드의 효용성을
기대하기 어려워졌기 때문에 필요하다면 창을 작게 만들어 한쪽에 띄워두는 정도로 사용하기 바란다.

\section{작업 흐름}

이 앱을 이용하는 작업 흐름은 
\begin{enumerate}[\quad 1)]\oblivoirlist
\item \ksemph{대응문자}를 입력하고
\item 필요한 문자를 필터와 메서드를 이용하여 생성한 후에
\item 클립보드를 통해 문서 에디터로 전송
\end{enumerate}
으로 진행하게 된다. \ksemph{대응문자}라고 한 것은 한글과 카나(仮名)를 \emph{로마자 알파벳}으로, 한자는 시스템 입력기를 통해서 직접 해당 문자를 입력한다는 의미이다.

대응문자와 한자의 입력과 문자 생성 방법에 대해서 다음 장부터 상세히 설명하려 한다.

\section{입력창과 입력 메서드}

다음 세 가지 단축키에 유의하자.
\begin{description}\oblivoirlist
\item [\hcrkey{Enter}] 현재까지 생성된 문자와 텍스트를 클립보드로 보내는 것이다.
\item [{\hcrkeys[+]{shift,enter}}] 입력 창 안에서 개행하는 것이다.
\item [\hcrkey{ESC}] 입력 창을 깨끗하게 비운다.
\end{description}
습관적으로 ``개행''을 위해 \hcrkey{Enter}를 치는 경우가 있을 수 있으므로 주의해야 한다.
만약 잘못 전송했거나 취소(\hcrkey{ESC})했다면, \hcrkeys[+]{Alt,Z} 단축키 또는 
\guielem{직전 전송 복구} 버튼을 눌러서 한 번에 한해서 취소하거나 전송한 것을 복원할 수 있다.

\begin{center}
\includegraphics[scale=.5]{scr03}
\end{center}

입력 메서드는 상단 가장 왼쪽의 콤보박스에서 선택하며 다음 다섯 가지 중에서 고를 수 있다.
\begin{description}\oblivoirlist
\item [\opt{pm}] 한글을 \pkg{pmhanguljamo}의 pm 입력방식으로 생성한다. 모든 음절을 성조 표지 또는 슬래시로 구분해주어야 한다.
\item [\opt{frkim}] 한글을 \pkg{pmhanguljamo}의 frkim 입력 방식으로 생성한다.
\item [\opt{mix}] 한글 자모의 대응문자는 pm 방식으로 하되, 음절 구성은 frk 방식을 따르는 절충 방식이다.
\item [\opt{rrk-ex}] 국가표준인 한글 로마자 표기법의 ``전자법''을 기본으로 20세기 초반 텍스트를 입력하기 위한 약간의 확장을 적용한 입력 방식이다. 옛한글을 완전하게 지원하지 않는다.
\item [\opt{Kana}] 로마자 입력으로부터 일본 문자를 얻어낸다.
\item [\opt{transparent}] 로마자 입력을 변환하지 않는다.
\end{description}
괄호 안에 약어가 있는데 이것은 이 앱의 \guielem{매직 태그}에서 사용하는 이름을 보인 것이다.

입력창에는 어떤 문자라도 올 수 있다. 다만 로마자 알파벳은 변환 대상이 되는 것이고 그 이외에는
문자를 있는 그대로 통과시킨다.

\section{한자 필터와 출력창}

생성된 문자와 텍스트가 출력창에 나타난다. 출력창을 디스플레이하는 폰트는 하단 폰트 선택 콤보박스에서
고를 수 있는데, 기본 폰트는 Noto Sans CJK KR이다. 폰트에 따라 자모 한글이나 일부 한자가 보이지 않을 수 있으므로 Noto Sans CJK KR을 쓰는 것을 권장한다.
\begin{center}
\includegraphics[scale=.5]{scr05}
\end{center}

입력창에 입력된 문자가 `한자'일 경우에는 \guielem{한자 필터}를
통해 변환된 한자를 생성한다.

\begin{center}
\includegraphics[scale=.5]{scr04}
\end{center}

\guielem{한자 필터}는 그림과 같이 네 가지가 있다.
\begin{description}\oblivoirlist
\item [없음] 한자를 변환하지 않고 입력된 그대로 출력한다.
\item [新字体] 현행 일본의 ``상용한자''에 쓰이는 글자체로 바꾸어서 출력한다.
\item [舊字體] 예전 일본의 한자 서체를 가리키는 말이기는 하나, 실제로는 대만식-한국식 한자(소위 繁體)와 거의 같다. 단지 서너 자의 자형이 다르다.
\item [简体] 중국어 간체자로 변환하여 출력한다.
\end{description}
이 각각을 지칭하는 약어는 \texttt{none}, \texttt{shin}, \texttt{kyu}, \texttt{sim}이다. 매직 태그에 활용하므로 잘 기억해두는 것이 좋다.

대응 한자의 선별은 \file{*.json} 사전 파일에 의한다. 이 사전은 필자가 상당한 노력을
들여 작성하기는 하였으나 완전하지 않을 것이다. 누락된 대응 한자가 있으면 알려주시기 바란다.


\section{정규화}

로마자 입력을 한글로 변환할 때에 기본적으로 각 입력값에 대응하는 \emph{자모 코드}가 출력된다.
NFD(Normalisation Form D)는 이른바 `첫가끝'이라 불리는 한글 표현 방식으로 자모 코드만으로
한글을 저장하는 것이고 NFC(Normalisation Form C)는 한글을 완성형 음절 코드로 저장하는 방식을 말한다.

즉, NFD를 선택하면 모든 한글이 전부 자모코드로 출력되지만 NFC를 선택하면 한글 음절 영역에 있는 11172자에 해당하는 것은 음절 코드로 전달된다(음절 영역에 존재하지 않는 한글, 즉 옛한글은 자모 코드 그대로 출력함).
현대 한글의 비중이 높다면 NFC를 선택해두는 것이 좋을 것이다.

\section{\LaTeX{} 확장 입력}

이 체크박스가 선택되어 있으면 다음 두 가지 기능이 활성화된다.\sidepar{v0.13}
\begin{itemize} \oblivoirlist
\item 확장 라틴 문자를 \LaTeX\ 입력 방식으로 얻어낸다. \verb|\"o|를 입력하면 \"o가 출력되는 식이다.
\item \verb|\x| 표기법으로 유니코드 코드포인트를 지정하여 문자를 입력할 수 있다. \verb|\x266B|는 ♫를 출력한다.
\end{itemize}

입력 메서드에 상관없이 위의 두 가지가 가능하지만 예컨대 \texttt{Dvořák}(\verb|Dvo\v{r}\'{a}k|)과 같이 아스키 문자가 섞인 
단어를 입력하려면 입력 메서드가 \opt{transparent}이거나 ``로마자 형태 보존'' 문법(p.~\pageref{sec:iden})을 사용하는 것이 좋을 것이다. 


\section{전송}

\hcrkey{Enter}를 누르거나 \guielem{전송} 버튼을 클릭하면, 
로마자 입력을 통해 구성한 출력을 클립보드로 보낸다. 편집 중인 문서로 돌아가서 붙여넣기(\hcrkeys[+]{Ctrl,V})하여 문서에 삽입하는 것이 이 앱의 목적을 실현하는 것이다.

\subsection{전송 후 초기화}

만약 \guielem{전송 후 초기화} 체크박스가 선택되어 있다면 
\hcrkey{Enter}를 눌러서 전송한 후 입력창은 깨끗하게 비워지고 새로운 입력을 받아들이게 한다.
반면 체크가 해제되어 있다면 전송한 후에도 입력 및 출력창의 텍스트와 문자가 그대로 남아 있게 된다.

\subsection{직전 창 전송}

이 체크박스가 선택되면 \hcrkey{Enter} 즉 \guielem{클립보드로 전송} 버튼은 
단순히 클립보드에 복사만 하는 것이 아니라 시스템 핫키를 조작해서 \hcrkeys[+]{Alt,Tab}과
\hcrkeys[+]{Ctrl,V}를 실행한다. Alt+Tab 키를 눌렀을 때 편집 중인 에디터의
편집 중인 위치에 커서가 있다면 \guielem{직전 창 전송}을 통해서 앱으로부터 에디터로
바로 텍스트를 전송하는 것 같은 효과를 얻을 수 있다.

다만, 시스템 핫키를 가로채는 것이라서 매번 성공하지 않을 수도 있고, \mbox{macOS}에서는
이 기능을 위하여 \textit{[개인정보 보호 및 보안]}의 \textit{[손쉬운 사용]} 권한을 얻어야
하는데, 설혹 실패하더라도 에디터 쪽에서 \hcrkeys[+]{Ctrl,V}하면 붙여넣기하는 것은
성공적으로 이루어질 것이다.

\chapter{공통 기호}

입력 메서드와 무관하게 특정한 기호 또는 스트링을 입력할 수 있도록 정의해두었다.

\section{문장부호}

문장부호는 \verb|^|(캐럿)과 함께 입력한다. 오늘날의 마침표, 쉼표, 물음표, 느낌표, 아스키 괄호 등은 키보드로부터 바로 입력한다. \tref{tab:caret}.

\begin{table}[htb]
\centering
\caption{캐럿과 함께 입력하는 문장부호}\label{tab:caret}
\begin{tblr}{
  width = \linewidth,
  colspec = {X[c,m] X[c,m] X[1.5,c,m] X[c,m] X[c,m] X[1.5,c,m]},
  hline{1,2,Z}={.8pt},
  vline{4}={},
  vline{2,3,5,6}={dotted},
  column{1,4}={font=\ttfamily},
  row{1} = {font=\rmfamily}
}
  입력 부호 & 출력 부호 & 비고 & 입력 부호 & 출력 부호 & 비고 \\
  \textcaret\~{} & ～ & 물결표 & \textcaret/ & / & 슬래시 \\
  \textcaret. & 。 & 고리점 & \textcaret= & 々 & 같은 글자 \\
  \textcaret, & 、 & 모점 & \textcaret: & ： & 전각 콜론 \\
  \textcaret> & · & 가운뎃점 & \textcaret; & ； & 전각 세미콜론 \\
  \textcaret< & ・ & 전각 가운뎃점 & \textcaret\textcaret & \textcaret & 캐럿 \\
  \textcaret- & ― & 긴 대시 & \textcaret\textbackslash & \textbackslash & 백슬래시 \\
  \textcaret * & … & 줄임표(절반) &  &  &  \\
\end{tblr}
\end{table}

\section{괄호류}

괄호류는 여닫는 부호이다. 괄호의 종류를 숫자로 나타내고 여닫는 부호 \verb|[|와 \verb|]|를
숫자 뒤에 붙여 괄호류를 입력한다. 여기서 괄호류는 키보드의 아스키 괄호가 아닌 CJK 부호들이다.

괄호의 종류를 나타내는 번호 0\~9는 \pkg{ksvqt} 패키지의 정의를 취한다. 그리고 \texttt{\#}과 함께 \#0\~{}\#5까지 여섯 개의 추가 기호를 정의하였다. \tref{tab:paren}.
따옴표는 괄호류로 취급한다. 특히 \verb|'|(small quote)를 문자 구성 부호의 일부로 쓰고 있는 \opt{frk}나 \opt{mix}에서 따옴표를 입력하려면 이 방법을 쓰는 것이 좋다.

\begin{table}
\centering
\caption{괄호류 입출력}\label{tab:paren}
\begin{tblr}{
  width = \linewidth,
  colspec = {X[c,m] X[c,m] X[c,m] X[c,m] X[c,m] X[c,m]},
  hline{1,2,Z}={.8pt},
  vline{4}={},
  vline{2,3,5,6}={dotted},
  column{1,2,4,5}={font=\ttfamily},
  row{1}={font=\rmfamily},
}
  여는 괄호 & 닫는 괄호 & 출력 & 여는 괄호 & 닫는 괄호 & 출력 \\
  0[ & 0] & 「 」 & 8[ & 8] & 〚 〛 \\
  1[ & 1] & 『 』 & 9[ & 9] & ｟ ｠ \\
  2[ & 2] & 〈 〉 & \#0[ & \#0] & ﹙ ﹚ \\
  3[ & 3] & 《 》 & \#1[ & \#1] & ﹛ ﹜ \\
  4[ & 4] & 〔 〕 & \#2[ & \#2] & ［ ］ \\
  5[ & 5] & 【 】 & \#3[ & \#3] & （ ） \\
  6[ & 6] & 〖 〗 & \#4[ & \#4] & ‘ ’ \\
  7[ & 7] & 〘 〙 & \#5[ & \#5] & “ ” \\
\end{tblr}
\end{table}


\section{\LaTeX\ 명령과 환경}

서두에 언급한 대로 우리는 \LaTeX{} 문서 작성을 염두에 두고 있다.
많지는 않겠지만 생성하고자 하는 텍스트에 \LaTeX{} 매크로가 포함되도록 하고 싶은 경우가 있다.
대표적인 것이 일본어에서 \cmdprint{\ruby}를 쓰는 경우일 것이다.

\pmime 는 백슬래시(\textbackslash)로 시작한 로마자 문자열이 공백(space), 중괄호(\{), 숫자를 만날 때까지를 별도로 보존하여 \LaTeX\ 매크로로 취급한다. 옵션 인자도 안전하게 보존되었다가 
복구된다. 첫 번째 중괄호로 주어지는 인자 안의 텍스트는 한글/카나로 변환한다.
또한 \verb|\begin|으로 시작하여 \verb|\end|로 끝나는 ``환경''의 이름은 한글로 변환하지 않는다. 

\begin{center}
\includegraphics[scale=.45]{scr06}
\end{center}

\noindent 주의할 것은 환경을 보존할 적에 관련있는 것은 \verb|\begin|과 그 다음 하나의 중괄호까지라는 것이다. 이 뒤에 옵션 인자 등을 덧붙이는 경우를 고려하지 않고 있다. 부득이한 경우에만 사용하는 것이 좋다.

\section{로마자 형태 보존}\label{sec:iden}

때로 외래어를 그 철자법 그대로 보여야 할 때가 있다. 이럴 경우를 위한 로마자 형태 보존 신택스를 마련해두었다. 즉, 백슬래시로 시작하고 단어의 끝에 백슬래시와 스페이스를 두는 것이다. 이 규칙은 단 하나의 단어에 대해서만 동작하며, 반드시 끝에 백슬래시와 스페이스가 존재해야 한다. 다음 예시를 보라.

\begin{center}
\includegraphics[scale=.45]{scr07}
\end{center}

만약 두 단어 이상이라면 첫 단어의 끝에 [백슬래시][스페이스]를 두어야 하므로 그 다음에 [스페이스]를 하나 더 입력해야 한다. 즉 예를 들면
\begin{verbatim*}
(\Good\  \morning\ )
\end{verbatim*}
과 같이 해야 한다. 여러 단어이고 길어질 때는 줄을 나누고 매직 태그를 이용하는 편이 더 낫다.

\section{매직 태그}

입력 메서드에 따라 입력 방법이 달라진다. 입력 메서드는 콤보 박스에서 선택할 수 있다.
한편, \pmime 는 행마다 입력 메서드와 한자 필터를 별도로 지정하여 강제하는 기능을 가지고 있다. 

행 머리에 \verb|%|와 메서드 지정 문자열, 그리고 한자 필터 지정 문자열을 적으면
그 행은 해당하는 입력 메서드와 한자 필터로 고정된다.
앞서 보인 다음 입력
\begin{verbatim}
mi'in'iranvn mardo dvdgi sirhji anihageni'oa i'inajvkei%
%jp (\ruby{許嫁}{iinazuke}),
eiqgeijimentv(\Engagement\ )ranvn mar'i isaqhagei gibbvgei dvrrinda.
\end{verbatim}
의 경우, 둘째 줄 맨 처음의 \verb|%jp|가 이 행의 텍스트는 일본어(카나)로 출력하라는 지시이다.

\begin{itemize}\oblivoirlist
\item \emph{입력 메서드 매직 태그}: \verb|%pm| (pm), \verb|%frk| (frkim), \verb|%mix| (mix), \verb|%rrk| (rrk-ex), \verb|%jp| (Kana), \verb|%clean| (transparent) 
이상 여섯 가지 매직 태그를 줄 수 있다.
\item \emph{한자 필터}: 한자 필터 태그는 반드시 메서드 태그 직후에 스페이스 없이 대괄호로 붙여야 한다. 즉 한자 필터만 지정할 수는 없고 항상 메서드 태그와 함께 지정해야 하는 것이다. 한자 필터 태그가 없으면 한자는 콤보 박스에 선택된 필터를 따른다.
\verb|[none]| (변환 없음), \verb|[shin]| (일본 신자체), \verb|[kyu]| (일본 구자체), \verb|[sim]| (간체)를 대괄호
안에 적어넣으면 된다. 단  입력 메서드 태그가 \verb|%clean|일 때에는 한자 태그가 아무 역할도 하지 않는다.
\end{itemize}

다음 예시 그림은 매직 태그를 활용하여 입력하는 예이다.

\begin{center}
\includegraphics[scale=.5]{scr08}
\end{center}

\chapter{한글}

\begin{figure}[htbp]
\centering
\includegraphics[width=\textwidth]{page32table}
\caption{한글 입력 대응 문자 요약}\label{fig:key}
\end{figure}

\fref{fig:key}\는 \pkg{pmhanguljamo} 패키지의 \href{https://texdoc.org/serve/pmhanguljamo-kdoc/0}{한국어판 설명서} p.32의 도표를 전재한 것이다. 
\pmime 는 기본적으로 이 입력 스킴을 그대로 따른다.

\section{pm, frkim 입력 메서드}

\subsection{pm}\label{sec:pm}

\opt{pm} 입력 방법은 \pkg{pmhanguljamo}에 고유한 방식이며, 다음과 같은 특징을 가진다.
\begin{enumerate}[\quad(1)]\oblivoirlist
\item 모든 한글 문자는 음절 단위로 `성조 부호'가 있는 것으로 간주한다. 즉 평성(\opt{/}), 상성(\opt{:}), 거성(\opt{;})의 성조 부호를 음절 끝에 \emph{반드시} 두어야 한다. 그 결과, 사성점이 사라진 이후의
표기법으로 적으려면 모든 문자가 ``평성''이라고 간주하여 음절 끝마다 슬래시를 쳐야 하게 되어 있다.
그 결과 한글 음절을 형성하기 위한 음절 분리 로직(오토마타)이 불필요하다.
\item 현재의 유니코드 한글 자모 코드 전체를 아우르도록 배정되어 있어서 이 입력 방법을 이용하면 ``모든 한글 문자''를 모두 구성할 수 있다. 여기서 ``모든 한글 문자''라 함은 초성, 중성, 종성 자모 코드를 이용하여 합성할 수 있는 수백만자를 가리키는 말이다. 
\item 자모를 단독자로 출력하기 위해서 유니코드 한글 호환자모 영역에 있는 글자는 대문자로 입력하여 얻을 수 있다. 만약 호환자모 영역에 정의되지 않은 단독자라면 초성-중성 채움문자를 이용해서 출력해야 할 것이다.
\item 특히 15세기의 중세 국어 표기를 조판하는 데는 이보다 좋은 방법이 현재로서는 없는 것 같다.
\end{enumerate}

\pkg{pmhanguljamo} 패키지 설명서의 예문 일부를 입력해보자.
%\begin{center}
%\includegraphics[scale=.4]{scr10}
%\end{center}
\begin{exampleonly}
na/ras;mar:ss@/mi; 中dyuf/國guig;ei; dar/a;
文mun/字jj@q;oa;ro; se/rv/ s@/m@s/di; a/ni;h@r/ss@i;

sai/mar/vr man/dvr/go sai/mar/vr ssv/nvn/ges/vn yu/haiq/i a/ni/ra yu/haiq/i/saq em/sug/ha/gei saiq/hoar/ei pir/yo/ha/ni/gga na/ta/na/nvn sa/sir/im/vr i/hai/hai/ya har/ges/i/da/.
\end{exampleonly}
\begin{quote}\exemfont
나랏〮말〯ᄊᆞ미〮 中듀ᇰ國귁〮에〮 달아〮
文문字ᄍᆞᆼ〮와〮로〮 서르 ᄉᆞᄆᆞᆺ디〮 아니〮ᄒᆞᆯᄊᆡ〮

새말을 만들고 새말을 쓰는것은 유행이 아니라 유행이상 엄숙하게 생활에 필요하니까 나타나는 사실임을 이해해야 할것이다.
\end{quote}

\subsection{frkim}

\opt{frkim}은 과거 1990년대의 \pkg{hangul\TeX}의 입력 방식을 현대적으로 재구성한 것이다.
그 연원이 명확하고 한글 대응 입력(transliteration, 轉字)의 이론적 기반이 있다.
대략의 특징은 다음과 같다.
\begin{enumerate}[\quad(1)]\oblivoirlist
\item 기본적으로 현대 한글을 기준으로 한다. 이 방식에서 핵심적인 음절 분리 로직은 ``모음 앞에는 단 한 글자의 자음만이 온다''는 것인데, 겹자음을 대문자로 쓰게 되어 있어서 이 로직을 적용하는 것이 가능하다.
\item 초성 `ㅇ'은 반드시 명시적으로\ksps{이 앱에서는 small quote를 이용하도록 했다} 밝혀주어야 한다. 이것을 생략하지 못하는 것이 원래의 \pkg{hangul\TeX}이었다.
\item 약간의 옛한글 자모를 출력할 수 있다.
\item 호환자모를 입력할 방법은 없다. 자모 단독자가 필요하다면 시스템 입력기에서 한글 단독자를 입력한다.
\end{enumerate}

\pkg{pmhanguljamo}의 \opt{frkim} 입력 방법은 여기에서 약간의 기능 확장이 이루어졌다.
\begin{enumerate}[\quad(1)]\oblivoirlist
\item 어절의 첫 초성 `ㅇ'을 생략할 수 있게 했다.
\item \verb|-|를 이용하는 음절 분리를 도입하여 복합 초성자를 식자할 수 있도록 했다. 그 결과
좀더 많은 옛한글을 입출력할 수 있게 되었다. 15세기 문헌을 100\% 처리할 수는 없지만 20세기 초반 문헌 정도는 자연스럽게 구현된다.
\end{enumerate}

\begin{exampleonly}
caig'yr hanguen gajigo iS'eSji'io. Gaman pioji'ei sonbadagmanhan caig'i'ji'io.
cesjaq'yr nemgimien nun'i nairigon haji'io.

nagpo'yi sdar'irnienni bando jinsaq oggieq gasda goaqhanjien'yisie jiegsoqja masna mijinjieqhoi h@'orca'yi
\end{exampleonly}
\begin{quote}\exemfont
책을 한권 가지고 있었지요. 까만 표지에 손바닥만한 책이지요.
첫장을 넘기면 눈이 내리곤 하지요.

낙포의 ᄯᅡᆯ일년니 반도 진상 옥경 갓다 광한젼의셔 젹송자 맛나 미진졍회 ᄒᆞ올차의
\end{quote}

\section{mix 입력 메서드}

원래 필자는 pm 메서드와 frkim 메서드를 내부적으로 숨기고 앱에서는 \opt{mix}만을 한글 입력 
방식으로 하려 했었다. pm 방식 최대의 장점인, ``모든 한글 자모를 모두 활용''하는 것을 살리고
그 대신 매번 슬래시를 쳐야 하는 불편을 해소하면서 성조점이 필요하다면 이것도 구현되도록 하려는 것이었다.

이를 위해서 (1) 자모 대응은 pm 방식의 것을 그대로 이용하되, (2) 음절 분리는 frkim 방식을
따르도록 하는 것, 이것이 \opt{mix} 메서드의 실체이다.
그러므로 \pmime\ 앱에서 한글 입력은 \opt{mix}를 기준으로 할 것을 권한다.

\opt{mix}에서 음절을 분리하는 방식에 대해 알아둘 사항이다.
\begin{enumerate}[\quad(1)]\oblivoirlist
\item frkim과 마찬가지로 어절(단어)의 첫머리가 아닌 한 초성 `ㅇ'은 \verb|'|로 밝혀야 한다.
\begin{verbatim}
nabogiga yeggye'ue
\end{verbatim}
`역겨워'에서 첫 `ㅇ'은 생략해도 된다. 그러나 이 단어의 세번째 음절 `워'는 \verb|'ue|로
초성 `ㅇ'을 밝혀야 한다.
\item \verb|-|(하이픈대시)는 새로운 음절이 시작하는 경계임을 나타내는 부호이다. 모음과 모음 사이에 너무 많은 자음이 와서 어디까지가 이전 음절의 종성이고 어디서부터가 다음 음절의 초성인지 알 수 없을 때 강제로 \verb|-|를 두어서 음절을 생성하게 한다. 특히 음절의 초성이 복합자음으로 이루어진 경우에는 반드시 그 시작을 명시해야 한다.
\begin{verbatim}
gasir-sdai'einvn
\end{verbatim}
``실ᄯᅢ'' 두 음절의 모음 `ㅣ'와 'ㅐ' 사이에 ㄹ, ㅅ, ㄷ이라는 세 자음이 들어 있다.
이것을 그냥 이어쓰면 ``싨대''가 될 수밖에 없다. 왜냐하면 종성 `ㄹㅅ'이 유니코드
테이블에 들어 있기 때문이다. 따라서 예시와 같이
명시적으로 \verb|-sd|라고 적어서 이것이 뒷 음절의 초성임을 밝혀야 한다.
\item 구현 과정에서 pm 방식의 슬래시(/)도 \verb|-|와 같은 역할을 하게 되었다. 그러므로 \verb|/|, \verb|'|, \verb|-| 세 부호가 음절 사이를 강제로 끊어주는 역할을 한다.
\item 호환 자모 입력 방법도 pm 방식과 같다.
\end{enumerate}

\pageref{sec:pm}페이지의 pm 입력 예제를 \opt{mix}로 다시 입력해보자. 초성 `ㅇ'을
밝히기 위해 \verb|'|를 쓴다는 사실에 유의한다.
\begin{exampleonly}
naras;mar:ss@mi; 中dyuf國guig;'ei; dar'a;
文mun字jj@q;'oa;ro; serv s@m@sdi; ani;h@rss@i;

saimar'vr mandvrgo saimar'vr ssvnvnges'vn yuhaiq'i anira yuhaiq'isaq emsughagei saiqhoar'ei pir'yohanigga natananvn sasir'im'vr ihaihai'ya harges'ida.
\end{exampleonly}
\begin{quote}\exemfont
나랏〮말〯ᄊᆞ미〮 中듀ᇰ國귁〮에〮 달아〮
文문字ᄍᆞᆼ〮와〮로〮 서르 ᄉᆞᄆᆞᆺ디〮 아니〮ᄒᆞᆯᄊᆡ〮

새말을 만들고 새말을 쓰는것은 유행이 아니라 유행이상 엄숙하게 생활에 필요하니까 나타나는 사실임을 이해해야 할것이다.
\end{quote}


\section{rrk 입력 메서드}

\opt{rrk}는 국가 표준 \<!한글 로마자 표기법>의 3장 8항 \<!전자법(轉字法)>에 의한
입력을 받아들인다. \<!한글 로마자 표기법> 자체는 음성적 표기라서 한글 입력 체계로는
쓸 수 없다. 예컨대 \texttt{songnisan}을 한글로 역전환하면 `\texttt{송니산}'밖에
더 얻겠는가? 즉 `소리'가 아니라 `표기'를 전자하는 예외 규정인 \<!전자법>에 의하여
입력해야 하는 것이다. 이 방식은 자모와 로마자 표기를 $1:1$로 대응시킨 체계인데
예컨대 초성 ㅇ은 \verb|-|이고 \verb|r|을 전혀 쓰지 않는 등의 특징이 있다.

이 방식은 원칙적으로 현대 한글밖에 적용할 수 없다. 왜냐하면 옛한글 자모에 대한 정의가
아예 없기 때문이다. 

\pmime 가 채택하고 있는 rrk는 \opt{rrk-ex}라고 해서 이 엄격한 제한을 조금 완화한
것이다. (1) 아래아(ㆍ)를 \opt{@}로 표기하게 한 것, (2) \verb|-| 부호의 의미를 확장해서
초성 `ㅇ'만이 아니라 음절 구분에도 적용하게 한 것, (3) 초성 ㄹ에 대해서 r도 허용하는 것 등이
완화된 점들이다.
그래서 ㅅ계 합용병서로 표기되는 초성 자음도 어떻게든 입력할 수 있게 하였다. 예컨대
\begin{verbatim}
%rrk g@r@mi p@r@ni sae deo-ug h@i-o
\end{verbatim}
이 입력은 
\begin{quote}
ᄀᆞᄅᆞ미 ᄑᆞᄅᆞ니 새 더욱 ᄒᆡ오
\end{quote}
라는 출력을 얻는다. 그러나 순경음 ㅂ(ㅸ)이나 그밖의 소실된 옛한글 자모를 입력할 방법이 없으므로
`20세기 초반 표기법' 정도까지만 적용될 따름이라고 보아야 한다. 

필자가 \opt{rrk}를 그다지 좋아하지 않는 이유는 여러 가지가 있으나, 결정적인 것은
`ㅓ'를 \opt{eo}에, `ㅡ'를 \opt{eu}에 할당한 모음 대응 방식이 마음에 들지 않아서였다.
불필요하게 타자수를 늘리고 음성적으로나 문자적으로나 식별이 어려워지는 이런 표기를 왜 굳이 해야
하는지 알 수가 없다. 매큔-라이샤워 방식이 사라지지 않는 이유도 여기에 있지 않을까?

대체로 말하여 \opt{rrk} 방식의 쓸모는 크지 않으나 이 방식을 굳이 포함한 이유는 어디선가 표준 로마자 표기법으로 입력된 텍스트를 역변환해야 할 가능성도 없지 않다고 보아서였다. 생각해보면 그런
표기들이 어차피 ``전자법''에 의한 것일 가능성이 대단히 희박하므로 이조차 그다지 효용이 있어
보이지는 않지만.

\<!님의 침묵> 처음 구절을 이 방식으로 입력해보겠다.

\begin{center}
\includegraphics[scale=.5]{scr09}
\end{center}


\section{네 가지 입력 메서드의 비교}


\begin{table}[htb]
\centering
\caption{입력 메서드의 지원 범위}\label{tab:methods}
\begin{tblr}{
  colspec = {c c c c c},
  hline{1-2,Z} = {},
  vline{2} = {dashed},
}
  입력 메서드 & 현대 한글 & 20세기초 한글 & 유니코드 (옛)한글 & 성조점 \\
  pm & \OK & \OK & \OK & \OK \\
  frk & \OK & \OK & \SEM & \SEM \\
  mix & \OK & \OK & \OK & \OK \\
  rrk & \OK & \SEM & \NON & \NON \\
\end{tblr}
\end{table}

\chapter{일본문자와 한자}

\pmime 는 일본어 입력기가 아니다. 그러나 일본어 문장을 알파벳 입력으로부터 얻을 수는 있다.

\section{히라가나와 카타카나}

히라가나는 소문자, 카타카나는 대문자로 입력한다. 표에서는 히라가나만 표시하였다.

\begin{table}[htbp]
\begin{widetext}
\hfill\begin{tblr}{
  width = .9\linewidth,
  colspec = {X[l,m] X[.6,c,m] X[l,m] X[.6,c,m] X[l,m] X[.6,c,m] X[l,m] X[.6,c,m] X[l,m] X[.6,c,m]},
  hline{1-Z}={dotted},
  hline{1,Z}={},
  hline{V} = {},
  vline{odd} = {},
  column{odd} = {font=\ttfamily}
}
  a & あ & i & い & u & う & e & え & o & お \\
  ka & か & ki & き & ku & く & ke & け & ko & こ \\
  ga & が & gi & ぎ & gu & ぐ & ge & げ & go & ご \\
  sa & さ & si, shi & し & su & す & se & せ & so & そ \\
  za & ざ & ji, zi & じ & zu & ず & ze & ぜ & zo & ぞ \\
  ta & た & ti, chi & ち & tu, tsu & つ & te & て & to & と \\
  da & だ & di & ぢ & du & づ & de & で & do & ど \\
  na & な & ni & に & nu & ぬ & ne & ね & no & の \\
  ha & は & hi & ひ & hu, fu & ふ & he & へ & ho & ほ \\
  ba & ば & bi & び & bu & ぶ & be & べ & bo & ぼ \\
  pa & ぱ & pi & ぴ & pu & ぷ & pe & ぺ & po & ぽ \\
  ma & ま & mi & み & mu & む & me & め & mo & も \\
  ya & や &  &  & yu & ゆ &  &  & yo & よ \\
  ra & ら & ri & り & ru & る & re & れ & ro & ろ \\
  wa & わ & wi & ゐ &  &  & we & ゑ & wo & を \\
  n' & ん &  &  &  &  &  &  &  &  \\
  fa & ふぁ & fi & ふぃ & fu & ふ & fe & ふぇ & fo & ふぉ \\
  va & ゔぁ & vi & ゔぃ & vu & ゔ & ve & ゔぇ & vo & ゔぉ \\
  xa & ぁ & xi & ぃ & xu & ぅ & xe & ぇ & xo & ぉ \\
  ?a & は & ?i & ひ & ?u & ふ & ?e & へ & ?o & ほ \\
\end{tblr}
\end{widetext}
\end{table}

엄격한 규칙을 따르기보다는 관행적 입력도 받아들이도록 중복 설계하였다.

\begin{table}[h]
\begin{tblr}{
  width = \linewidth,
  colspec = {X[l,m] X[.6,c,m] X[l,m] X[.6,c,m] X[l,m] X[.6,c,m]},
  hline{1-Z} = {dotted},
  hline{1,Z} = {},
  vline{odd} = {},
  column{odd} = {font=\ttfamily},
}
  kya & きゃ & kyu & きゅ & kyo & きょ \\
  gya & ぎゃ & gyu & ぎゅ & gyo & ぎょ \\
  sha, sya & しゃ & shu, syu & しゅ & sho, syo & しょ \\
  ja, zya, jya & じゃ & ju, zyu, jyu & じゅ & jo, zyo, jyo & じょ \\
  cha, tya & ちゃ & chu, tyu & ちゅ & cho, tyo & ちょ \\
  nya & にゃ & nyu & にゅ & nyo & にょ \\
  hya & ひゃ & hyu & ひゅ & hyo & ひょ \\
  bya & びゃ & byu & びゅ & byo & びょ \\
  pya & ぴゃ & pyu & ぴゅ & pyo & ぴょ \\
  mya & みゃ & myu & みゅ & myo & みょ \\
  rya & りゃ & ryu & りゅ & ryo & りょ \\
  xya, \~{}a & ゃ & xyu, \~{}u & ゅ & xyo, \~{}o & ょ \\
\end{tblr}
\end{table}

발음 ん는 단독자로 입력할 때는 \verb|n'|이다. 입력 문자열의 문맥에 따라 \verb|n|을
ん로 바꾼다. 그러므로 \verb|konnichi?a|가 ``こんにちは''로 자연스럽게 변환한다.

\verb|?a|와 \verb|?e|는 `は'로 적고 `wa'로 읽는 조사, 그리고 `へ'로 적고 `e'로
읽는 조사를 위한 입력 방식이다. 음성 지각과 입력 문자의 상위를 줄이고자 하는 의도에서
정의하였다.

촉음 っ를 단독자로 입력할 때는 \verb|xt|로 한다. 단어 안에서는 굳이 매번 \verb|xt|라고 하지 않아도 입력 문자열에서 중복 자음을 
프로그램이 자연스럽게 잘 처리한다. \verb|chotto|ちょっと \verb|itterasshai|いってらっしゃい.

히라가나는 별도의 장음 부호가 없고 문자 그대로 입력한다. \verb|toukyou|とうきょう \verb|oosaka|おおさか. 카타카나의 장음 부호는 \verb|--|로 입력한다. \verb|KO--HI--|コーヒー \verb|BI--RU|ビール.

`코가키가나'(小書き仮名)의 하나인 ヶ는 \verb|xk|로 입력한다. 이 글자는 주로 `ga' 또는 `ka'로 읽을 때가 많다. 3ヶ月 関ヶ原.

\section{일본어 문장 부호}

Kana 입력 메서드에는 공통 부호 이외에 이 입력 방식에만 적용되는 문장부호가 있다.

\verb|{noma}|와 \verb|{onaji}|는 각각 々, ゝ를 출력한다. 둘 다 앞 글자의 반복을 나타낸다. 보기: 日々 人々. 
이 가운데 \verb|{noma}|는 공통 문장부호 \verb|^=|로 입력해도 같은 결과를 얻는다. 
시마자키 토손의 \<!첫사랑> 중에 다음과 같은 표기가 나온다.

\verb|wagako{onaji}ronaki tameikino| わがこゝろなきためいきの

\begin{figure}[ht]
\includegraphics[scale=.45,center]{scr11}
\end{figure}

Kana 입력 메서드에서는 마침표와 쉼표, 그리고 \verb|`|(backtick), \verb|'|(single quote) 부호를 그림과 같이 출력한다. 마지막 단어 「えん」을 \verb|`en''|으로 입력했는데 마지막에 싱글 따옴표가 두 개인 이유는 \verb|n'|이 단독으로 ん을 의미하기 때문이다.

일본어 입력 상태라도 스페이스 즉 띄어쓰기에는 아무런 조작을 가하지 않는다. 띄어쓰는 방식으로 입력한 텍스트를 띄어쓰기 없이 출력하는 것은 \LaTeX 에게 맡겨져 있다. 일본어 문장을 식자할 때 어떤 패키지를 쓰느냐에 따라 이를 조절할 방법이 있을 것이다.

\section{칸지}

우리말 문장은 이제 한자 없이 쓰는 것이 거의 표준이 되었다. 그러나 한자 없는 일본어 문장이란
상상하기 힘들다. 카나만이 아니라 한자도 섞인 일본어 `문장'을 얻어내려면 어떻게든 이 한자라는
문제를 극복해야 한다.

일본어 입력기의 핵심적 기능은 한자 변환이다. \pmime 가 ``일본어 입력기가 아니다''고 말한
이유는 이런 기능을 탑재하지 않았기 때문이다. 

이 앱의 목적을 다시 생각해보면 우리는 한국어-한글 환경에서 문서를 작성하고 있는 것이다.
그런데 인용이든 다른 이유에서든 일본어 단어와 문장을 문서에 삽입해야 하는 상황이다.

\subsection{한국어 입력기로 입력한 한자}

운영체제에 딸린 한국어 입력기는 한자 입력 기능을 갖추고 있다. 한국어 입력기를 통하여
입력하기 \emph{쉬운} 글자는 한국식 한자이다. 

한국어 입력기로 한국식 한자를 입력하더라도 출력을 일본의 신자체 한자로 바꾸자는 것이
\pmime 의 처음 디자인 컨셉트였다. 이것이 \guielem{한자 필터}로 구체화된 것이다.
일본어 원어 사용자라면 상상하지 않을, 철저하게 한국어 사용자의 편의적 입장에서 만든 것임을 알 수 있겠다.

한국어 입력기의 한자 입력 방법으로 다음과 같이 입력하였다고 하자.

\centerline{沸點 化學 液體 氣體 溫度}

\guielem{한자 필터} \<!\opt{新字体}>를 거쳐 출력은 이렇게 나타난다.

\centerline{沸点 化学 液体 気体 温度}

그림의 한자는 모두 한글 입력기로부터 입력한 것이다. 예컨대 \<2溫>은 비슷하게 생겼지만 다른 글자 \<2温>으로 잘 바뀌어 있다.

\begin{center}
\includegraphics[scale=.5,center]{scr12}
\end{center}

한국식 한자와 일본 신자체 변환 정보는 \file{shinjitai.json} 사전 파일에 등록되어 있다.
현재 거의 모든 상용한자 범위 신자체에 대한 항목이 망라되어 있기는 하나, 이따금 이체자 등의 이슈로 
원하는 결과가 나오지 않을 때가 없지 않을 것이라 생각한다. 이럴 경우에는 해당 사전 파일을 수정하여
쓰시고, 꼭 제작자에게 알려주시기 바라는 바이다.

\subsection{칸지 입력 보조창}

그런데, 한국어 한자 입력기만으로는 입력이 불편한 일본식 한자들이 있다. 소위 `코쿠지'(国字)라는 것이나 일상 한자라도 한국어 한자 입력기에서 너무 뒷 순번에서 찾아야 하는 것도 있다.
이런 한자들의 입력 효율을 높이기 위해서, 다음과 같은 한자 입력 태그를 도입하였다.
\begin{center}
\includegraphics[scale=.5]{scr13}
\end{center}

이 \verb|{h:| 입력 태그에 의한 한자는 \file{customkanji.json}에 정의되어 있는데 
입력 키가 어떻게 정의되어 있는지 알기 어려울 것이라서 이 파일의 내용을 보여주고 검색할 수 있게
하는 \guielem{한자 입력 보조창}을 \hcrkey{F10}으로 열게 하였다.

\begin{description}\oblivoirlist
\item [\hcrkey{F10}] 한자 입력 보조창을 열거나 닫는다.
\item [{\hcrkeys[+]{Ctrl,F}}] 한자 입력 보조창의 검색 필드로 커서를 이동한다.
\item [\hcrkey{Enter}] 검색 필드에서 숫자 또는 \hcrkey{Enter}를 누르면 입력창에 해당 한자의 키값이 들어가고 출력창에 한자가 나타나며 포커스가 입력창으로 이동한다. 만약 한자 검색을 계속해야 한다면 \hcrkeys[+]{Ctrl,F}로 한자 검색 위치로 포커스를 옮긴다.
\end{description}

\hcrkey{ESC}의 동작에 대해 주의하여야 한다. 이 키는 입력창의 내용을 비우는 역할을 하는데, 포커스가 한자 입력 보조창에 있을 때만 보조창을 닫는 역할을 한다. 그러나 포커스를 일일이 확인하지 않으면 실수로 (물론 복구할 수 있지만) 입력창 내용을 지우게 될지 모르니 보조창의 여닫기는 \hcrkey{F10}으로만 하도록 하는 것이 좋다.

\begin{center}
\includegraphics[scale=.5]{scr14}
\end{center}

그림에서 보듯이 사전에 등록되지 않은 많은 한자들이 있다. 이 사전은 편의를 위해서 만든 것이고
누차 말하듯이 이 프로그램이 일본어 입력기로 설계한 것이 아니기 때문에 완전한 일본어 입력 시스템에
준하는 정도의 내용을 갖추고 있지 않다. 개구리 蛙자는 사전에 등재되지 않은 것이라 한국어 입력기로부터
직접 입력하였다.
\begin{quote}
古池や、 \ruby{蛙}{かわず} 飛び込む 水の音
\end{quote}

입력 보조창의 찾기 필드에 입력하는 텍스트에 대해 점진적 일치로 찾은 결과를 보여준다. 예컨대 \verb|aka|를 입력하면 
\opt{aka}, \opt{baka}, \opt{naka} 등의 일치하는 단어가 뜨게 되는데, \verb|aka|로 시작하는 단어만 검색하려면
\verb|:|부터 입력한다. \verb|:aka|. 그러면 예컨대 \opt{h:aka}, \opt{h:akari} 등이 나타날 것이다.

\paragraph{새 단어 등록}
사용자는 \file{customkanji.json}의 내용을 자신에게 맞게 보충해가면서 쓸 수 있다.
\guielem{새 단어 등록}을 누르고 검색 키값과 한자를 차례로 입력하여 등록하면 그때부터
해당 키값으로 한자를 불러낼 수 있다.

이처럼 개인화한 \file{customkanji.json}은 별도로 백업해두는 것이 좋겠다. Windows는
설치 폴더에, 그리고 macOS에서는 \path{~/.config/pmimegui/}에 있는 파일이 새 단어
등록 대상이다.

\begin{center}
\includegraphics[scale=.45]{scr15-1}
\includegraphics[scale=.45]{scr15-2}
\end{center}

비록 계속해서 ``이 프로그램은 일본어 입력기가 아니다''고 강조하고 있지만
실은 간이 일본어 입력기처럼 쓰는 것이 아주 불가능하지는 않은 정도가 되었다고 생각하고 있다.

\subsection{일본어 신자체에 대응하는 한국 한자}

매직 태그의 한자 태그를 \opt{[han]}으로 하면 입력 텍스트의 일본어 신자체를 한국식
한자로 필터링한다. 콤보 박스에서는 선택할 수 없고 매직태그로만 쓰도록 하고 있는 기능이다.
그 이유는 이 기능이 별도의 사전을 참조하는 것이 아니라 
한국식 한자를 일본어 신자체에 대응시킨 사전을 역변환하여 보여주는 것일 뿐이라서
때로 여러 자가 하나의 자체에 대응하기도 하고 일부 당연한 변환이 결락되는 등
그 자체로 신뢰할 기능으로 동작하는 것은 아니기 때문이다. 한 예를 들자면 `よゆう'로 읽는
`余裕'에 이 필터링을 적용하면 `\{餘|余\}裕'라는 결과를 얻는데 이것은 余에 대응하는
글자가 여럿이라는 의미이고, 이런 식으로 $1:\mathrm N$ 대응의 결과 원치 않는 변환이
일어나는 경우도 있다. 그것은 한국어의 한자어와 일본어의 한자어가 일치하지 않기 때문이다.
이 문제를 해결하려면 단어 수준의 변환 사전을 만들어야 할 것이나, 우리 목적에
부합하지 않고 너무 방대한 작업을 요구한다.

그러므로 이 기능은 오직 단순 참고용으로만 활용해야 할 것이다. 예컨대 다음 예문에서
国, 号같은 한자가 國, 號로 바뀌어 출력되는 것을 확인하는 정도이다.

\begin{exampleonly}
%jp[han] 国境の長いトンネルを抜けると雪国であった。夜の底が白くなった。信号所に汽車が止まった。
\end{exampleonly}
\begin{quote}
國境の長いトンネルを拔けると雪國であった。 夜の底が白くなった。 信號所に汽車が止まった。
\end{quote}

\section{중국어 간체}

중국어에 대해서는 별도의 입력 방법을 갖추고 있지 않다. 단지 한국식 한자로 입력한 문자를
중국어 간자체로 필터링하는 기능만이 있을 뿐이다. 번체를 간체로 바꾸어주는 간략화한 버전이라고
보면 된다. 

따라서 중국어 입력기를 대용할 수 없다.
그래도 간자체를 생성할 필요가 있을 때 이따금 쓸모가 있기를 바랄 뿐이다. 번체$\;\Leftrightarrow\;$간체 변환 도구는 너무나 많다.
중국어 입력기 비슷하게 동작하게 하려면 적어도 pinyin(拼音)으로부터 한자를 얻어내는
기능이 있어야 할 것이나 우리의 고려사항이 아니었다.
일본어에 대해서는 대응하고 있는 `역변환'도 간체에 대해서는 적용되지 않는다. 

\begin{center}
\includegraphics[scale=.5]{scr16}
\end{center}

\chapter{커맨드라인 인터페이스}\label{sec:cli}

\pmime 는 GUI로 동작하지만 CLI 인터페이스로 사용할 수 있다.
특히 대량의 자료를 파이프나 파일로 처리할 때 유용하다.

\begin{center}
\includegraphics[width=\linewidth]{scr17}
\end{center}

\section{커맨드라인 실행}

\subsection{Windows}

인스톨러로 설치한 경우에는 \util{pmimec}라는 명령행 유틸리티를 터미널에서 실행한다.

단독 실행 파일과 소스를 다운로드한 때는 \file{pmimegui.py} 스크립트를 Python으로
실행하면서 인자를 부여하면 된다. CLI 실행시에는 PyQt6 같은 GUI용 패키지를 요구하지
않으므로 기본 설치된 Python으로도 동작한다.

\subsection{macOS}

\file{dmg}로 \path{/Applications}에 app 번들을 설치하였을 때는 다음 명령을
터미널에서 실행한다.
\begin{verbatim}
/Applications/pm\ IME\ GUI.app/Contents/MacOS/pm\ IME\ GUI --help
\end{verbatim}
필요하다면 간단한 셸 스크립트를 만들어서 실행해도 좋을 것이다.
\begin{exampleonly}
#!/bin/sh
/Applications/pm\ IME\ GUI.app/Contents/MacOS/pm\ IME\ GUI $*
\end{exampleonly}

\subsection{Python Script}

터미널에서 직접 실행하면 된다.
\begin{verbatim}
python <path/to>/pyimegui.py --help
\end{verbatim}

\section{인자}

\subsection*{도움말과 버전}
\begin{description}\oblivoirlist
\item [\opt{-v}] \verb|--version|. 버전 번호를 보여준다.
\item [\opt{-h}] \verb|--help|. 도움말을 보여준다.
\end{description}

\subsection*{메서드와 필터}
\begin{description}\oblivoirlist
\item [\opt{--mode}] 변환 메서드 지정. \opt{pm}, \opt{frk}, \opt{mix}, \opt{rrk}, \opt{jp}, \opt{clean} 중에서 한 가지를 부여한다.
\item [\opt{--hanja}] 한자 필터. 기본값은 \opt{none}이고, \opt{shin}, \opt{kyu}, \opt{sim}, \opt{han} 중에서 한 가지를 부여한다.
\end{description}

\subsection*{보조 동작}
\begin{description}\oblivoirlist
\item [\opt{-n}] \verb|--nfc|. 출력 한글을 NFC로 정규화. 이 옵션을 주지 않으면 NFD.
\item [\opt{-c}] \verb|--clip|. 출력 스트링을 클립보드로 복사.
\end{description}

\subsection*{파일}
\begin{description}\oblivoirlist
\item [\opt{-i}] \verb|--input|. 읽을 파일 이름을 지정한다.
\item [\opt{-o}] \verb|--output|. 출력을 콘솔 대신 파일로 저장할 때 파일 이름을 부여한다.
\item [\opt{--macro}] 입력 파일에서 변환할 부분에 해당하는 명령어를 지정한다. 예를 들면 \cmdprint{\jamoword}의 인자만을 변환하려 한다면 \verb|--macro \jamoword|.
\item [\opt{--env}] 입력 파일에서 변환할 대상에 해당하는 환경의 이름을 지정한다. 예를 들면 \env{jamotext} 내부의 것만을 변환하려 할 때 \verb|--env jamotext|.
\item [\opt{-r}] \verb|--raw|. 파일을 처리할 때에 파일 전체를 한글 텍스트로 간주한다. 즉 macro나 env 없이 파일 전체를 변환한다.
\item [\opt{--keep}] 이 옵션이 주어지면 예컨대 \verb|\jamoword{han/gvr}|을 ``한글''로 변환하는 것이 아니라 \verb|\jamoword{한글}|로 변환한다. 즉 \LaTeX\ 명령어 형태를 그대로 남긴다. 환경에 대해서도 같다. 
\end{description}
\verb|--macro|와 \verb|--env|는 각각 하나씩밖에 지정할 수 없지만 매크로 하나, 환경 하나를 별도로 지정하는 것은 가능하다. 예: \verb|--macro \hg --env hangul|. 이 경우 \cmdprint{\hg}의 인자와 \opt{hangul} 환경 내부의 텍스트를 한글로 변환하려 할 것이다.

\section{사용례}

실행 명령은 \util{pmimec}인 것으로 서술한다.

\subsection{변환 텍스트를 인자로 부여}

변환할 입력 문자열을 따옴표(\verb|"|)로 묶어 전달한다.
\begin{exampleonly}
$ pmimec --mode mix -n -c "hataq'ei gase bar ssisgo, juqtaq'ei gase son ssisgo, saqtaq'ei gase nac ssisgo onera"
\end{exampleonly}
\begin{quote}\exemfont
하탕에 가서 발 씻고, 중탕에 가서 손 씻고, 상탕에 가서 낯 씻고 오너라
\end{quote}

\begin{exampleonly}
$ pmimec "幼時kara 父?a, {h:watashi}niyoku, 金閣nokotowo 語tta." --mode jp --hanja shin -c
\end{exampleonly}
\begin{quote}\exemfont
幼時から 父は、 私によく、 金閣のことを 語った。
\end{quote}

위 두 경우 모두 \verb|-c| 옵션을 주어서 클립보드에 복사하였으므로 실행 후에 다른
아무 곳에서나 \hcrkeys[+]{Ctrl,V}로 결과를 붙여넣을 수 있다.

\subsection{파일로부터 입력}

아무런 옵션도 주지 않았을 때 입력되는 파일의 내용 중에 \cmdprint{\jamoword}의
인자가 있는지를 찾아서 그 인자 안의 텍스트가 pm 입력이라고 간주하고 변환을 시도한다.
그러므로 한글 부분을 \cmdprint{\jamoword}로 pm 입력해둔 파일이 있다면 단순히
\begin{verbatim}
pmimec -i inputfile.txt 
\end{verbatim}
이렇게만 해도 \cmdprint{\jamoword}의 인자가 한글로 변환되어서 나올 것이다.

\cmdprint{\jamoword} 대신 \cmdprint{\hg}를 써보자.
\begin{figure}[htbp]
\begin{boxedverbatim}
Jikji (Korean: \hg{jig/ji/sim/cei/yo/jer}) is the abbreviated
title of a Korean Buddhist document whose title translates 
as "Anthology of Great Buddhist Priests' Zen Teachings".
\end{boxedverbatim}
\caption{입력 파일의 내용}\label{fig:ec}
\end{figure}

\fref{fig:ec}\과 같은 내용을 가진 \file{input.txt}가 있다고 하자.
이 파일의 내용 중에서 \verb|\hg|의 인자로 주어진 로마자는 한글로 변환해야 한다.

\begin{exampleonly}
$ pmimec -i input.txt --mode pm -n --macro \hg
\end{exampleonly}
\begin{quote}\exemfont
Jikji (Korean: 직지심체요절) is the abbreviated
title of a Korean Buddhist document whose title translates
as "Anthology of Great Buddhist Priests' Zen Teachings".
\end{quote}

만약 \cmdprint{\hg}라는 매크로를 보존해야 한다면

\begin{exampleonly}
$ pmimec -i input.txt --mode pm -n --macro \hg --keep
\end{exampleonly}
\begin{quote}\exemfont
Jikji (Korean: \cmdprint{\hg}\{직지심체요절\}) is the abbreviated
title of a Korean Buddhist document whose title translates
as "Anthology of Great Buddhist Priests' Zen Teachings".
\end{quote}

환경(environment)에 대해서도 같은 방식으로 할 수 있다. 환경일 때 \verb|--keep|을
부여하면 \verb|\begin{...}|와 \verb|\end{...}|이 그대로 보존된다.

결과를 파일로 저장하고자 하면 \verb|-o| 옵션과 파일이름을 준다.

만약 파일의 일부\ksps{매크로 또는 환경}가 아니라 파일 전체가 통째로 한글에 해당하는
변환 대상이라면 \verb|-r| 옵션을 주어서 실행한다.

\subsection{파이프}

표준 입력을 통해 파이프 처리할 수 있다. 파이프를 통해 넘어오는 입력 문자열에 대하여
\cmdprint{\jamoword}를 찾는 것이 기본값이므로 입력 문자열 전체를 변환 대상으로 하려면
\verb|-r| 옵션을 반드시 주어야 한다.

\begin{exampleonly}
$ echo "jikji \jamoword{jig/ji}" | pmimec --mode pm
\end{exampleonly}
\begin{quote}\exemfont
jikji 직지
\end{quote}

\begin{exampleonly}
$ echo "jug'edo ani nunmurhvrri'urida" | pmimec --mode mix -r
\end{exampleonly}
\begin{quote}\exemfont
죽어도 아니 눈물흘리우리다
\end{quote}

\chapter{나가며}

이 지침서에서 다룬 내용을 간단히 정리하며 마무리하고자 한다.

《pm IME GUI》는 시스템 입력기를 전환하지 않고도 옛한글, 일본 문자, 한자, 중국어 간체 등을 하나의 창에서 처리할 수 있도록 만든 보조 도구이다. 애초의 동기는 지극히 개인적인 것이었다. 필자가 다루고 있던 1920~1930년대 문헌은 오늘날의 한글 표기법과도, 15세기 이래의 옛한글과도 다른 독특한 표기 관습을 가지고 있었고, 이를 입력하기 위해 매번 시스템 입력기를 전환하는 일은 집필의 흐름을 자꾸만 끊어놓았다. 이 불편을 해소하기 위해 만든 도구가 여기까지 이르렀다.

이 앱의 핵심 메커니즘은 필자가 처음부터 새로 고안한 것이 아니라 pmhanguljamo 패키지의 자모 입력 방식을 차용한 것이다. 한글 자모를 로마자로 입력받아 유니코드 한글 자모 코드로 합성해내는 이 방식이 없었다면 《pm IME GUI》도 존재하지 않았을 것이다. pm과 frkim이라는 두 입력 방식의 이름 자체가 이 패키지에서 유래한 것임은 5장에서 이미 설명한 바와 같다. mix 방식은 이 두 방식의 장점을 절충하려는 시도였고, rrk 방식은 국가 표준 로마자 표기법과의 호환을 위해 추가로 구현해본 것이다.

이 프로그램이 다루는 범위는 애초에 넓지 않다. 한글, 일본 문자, 그리고 약간의 한자 변환 정도가 전부이며, 일본어 입력기나 중국어 입력기를 표방하지도 않는다. 6장에서 여러 차례 밝혔듯 한자 변환에 쓰이는 사전 파일들은 필자 개인의 작업으로 채워진 것이라 누락과 오류가 있을 수밖에 없다. 신자체·구자체·간체 변환이 어색하게 나오는 경우를 발견하면 언제든 알려주시기 바라며, customkanji.json과 shinjitai.json은 애초부터 사용자가 직접 고쳐 쓰도록 열어둔 파일이니 필요에 따라 자유롭게 보완해서 사용하시면 된다.

돌이켜보면 이 도구는 처음부터 거창한 목표를 가지고 시작된 것이 아니었다. 단지 스스로의 작업 흐름에서 반복되던 사소하지만 성가신 불편 하나를 해결하고자 했을 뿐이다. 그러나 비슷한 종류의 텍스트, 즉 근대 초기의 표기법이나 다국어가 섞인 인용문을 다루어야 하는 다른 사용자에게도 이 도구가 조금이나마 도움이 될 수 있으리라는 생각에 공개하게 되었다.

《pm IME GUI》는 앞으로도 필자의 실제 작업 과정에서 필요에 따라 조금씩 손보아 나갈 예정이다. 완성된 도구라기보다는 계속 다듬어지는 도구에 가깝다고 보는 편이 정확할 것이다. 버그를 발견하거나 누락된 한자, 지원되지 않는 표기가 있다면 언제든 알려주시기 바라며, 이 지침서 역시 프로그램의 변화에 맞추어 계속 갱신해나갈 것이다.

부족한 점이 적지 않은 도구이지만, 이 지침서를 끝까지 읽어주신 독자들께 감사의 말씀을 전한다.


\section*{이런 분들도 써보시기 바랍니다}

이 지침서는 지금까지 다소 전문적인 용도, 즉 20세기 초반 표기법이나 옛한글 문헌을 본격적으로 조판하는 상황을 중심으로 설명하였다. 그러나 《pm IME GUI》는 그보다 훨씬 가벼운 용도로도 유용하다.

예를 들어 한글 워드프로세서의 역사에 관한 글을 쓴다고 하자. 본문 대부분은 현대 한글이지만, 초기 워드프로세서가 자모 조합에 어려움을 겪었다는 사실을 설명하려면 "ᄒᆞᆫ글"이라는 표기 한 번쯤은 등장할 수밖에 없다. 이 한 글자를 위해 옛한글 입력기를 새로 설치하거나 시스템 입력기를 전환하는 것은 명백히 배보다 배꼽이 큰 일이다. 이럴 때는 mix나 pm 모드로 그 단어만 잠깐 입력하고 클립보드로 보내는 편이 훨씬 간편하다.

일본어도 마찬가지다. 일본어를 전혀 할 줄 몰라도 검색엔진에 일본어 단어 하나를 넣어야 할 때, 또는 인용문에 짧은 일본어 문장을 넣어야 할 때가 있다. 이런 경우에도 일본어 IME를 평소에 켜둘 필요 없이, Kana 모드로 로마자를 그대로 쳐서 필요한 순간에만 히라가나·가타카나를 얻어내면 그만이다.

즉 이 앱의 진짜 쓸모는 두 가지 상반된 상황 모두에서 나온다.
\begin{itemize}\oblivoirlist
\item 옛한글이나 일본 문자, 여러 자체의 한자를 \emph{본격적으로}, \emph{대량으로} 다루어야 하는 경우 (이 경우 5장\~7장에서 설명하는 각 입력 메서드와 매직 태그, CLI를 익혀두면 좋다)
\item 평소에는 전혀 쓸 일 없는 문자를 \emph{아주 가끔}, \emph{한두 글자만} 넣어야 하는 경우 (이 경우 5장의 표 하나, 또는 6장의 가나 대응표 하나만 봐도 충분하다)
\end{itemize}

후자에 해당하는 독자라면 이 지침서를 처음부터 끝까지 읽을 필요는 없다. 3장에서 기본 조작(입력창--출력창--클립보드 전송)을 익히고, 필요한 문자 체계에 해당하는 대응표(5장 그림 1, 6장 히라가나 표)만 옆에 띄워두고 참고하면 그것으로 충분하다. 시스템 입력기를 켰다 껐다 하는 번거로움 없이, 필요한 순간에만 필요한 문자를 만들어 쓰는 것, 이것이 이 도구가 애초에 추구한 목적이다.


\end{document}
